\section{Methodology}
\label{sec:method}

\subsection{Design principles}
Two rules shape the design. (1) \emph{The LLM never sees the raw image}, only structured, numbered findings.
This shrinks the attack surface and makes every statement in the report traceable to a finding id. (2)
\emph{Facts come from code, not from the LLM.} Cipher confidence, key strength, recovered keys and the risk
score are computed deterministically; the LLM only writes them up, and a validator checks that it did so
faithfully.

\subsection{System architecture}
Fig.~\ref{fig:block} shows the pipeline. Five analysis stages write findings to SQLite; a fusion stage turns
them into scores; the report stage produces either a rule-based report (always available) or a defended LLM
report. Each stage is independent, so a missing input (no symbol table, no capture, no API key) skips that
stage instead of failing the run. The package has about 4{,}000 lines of Python in 22 modules; a
standard-library web server provides the dashboard, and a static build of it is published on GitHub Pages.

\begin{figure*}[t]
\centering
\resizebox{\textwidth}{!}{%
\begin{tikzpicture}[stage/.style={blk, text width=66mm, minimum height=9mm}]
  \node[stage] (s1) at (6.2,0)     {\textbf{1 Memory}: Volatility~3 (Windows, Linux): processes, command lines, sockets, malfind, bash history, USB; region grading};
  \node[stage] (s2) at (6.2,-1.35) {\textbf{2 Static}: strings, PE headers, imports, section entropy, YARA on dumped regions};
  \node[stage, fill=green!12] (s3) at (6.2,-2.7) {\textbf{3 Crypto (core)}: entropy, cipher ID (Dempster-Shafer), RSA tests (NIST), AES / ChaCha20 / RSA key recovery};
  \node[stage] (s4) at (6.2,-4.05) {\textbf{4 Credentials}: JWT, OAuth/API tokens, OTP secrets (redacted)};
  \node[stage] (s5) at (6.2,-5.4)  {\textbf{5 Network}: pcap via tshark: TLS version and suite, DNS, SNI};
  \node[io, text width=17mm] (dump) at (0,-2.7) {Memory\\image};
  \node[io, text width=17mm] (pcap) at (0,-5.4) {pcap\\(optional)};
  \draw (dump.east) -- (1.8,-2.7);
  \draw (1.8,0) -- (1.8,-4.05);
  \draw[->] (1.8,0) -- (s1.west);
  \draw[->] (1.8,-2.7) -- (s3.west);
  \draw[->] (1.8,-4.05) -- (s4.west);
  \draw[->] (pcap.east) -- (s5.west);
  \draw[->] (s1.south) -- node[right, font=\scriptsize]{dumped regions} (s2.north);
  \foreach \s in {s1,s2,s3,s4,s5} \draw (\s.east) -- (\s.east -| 10.0,0);
  \draw (10.0,0) -- (10.0,-5.4);
  \node[db] (db) at (11.3,-2.7) {SQLite\\findings\\F1, F2, ...};
  \draw[->] (10.0,-2.7) -- (db.west);
  \node[blk, text width=22mm] (fuse) at (14.0,-1.2) {Fusion: risk,\\range, ATT\&CK,\\campaign links};
  \node[blk, text width=20mm] (rule) at (14.0,-4.2) {Rule-based\\report};
  \node[llm, text width=33mm] (llm)  at (17.8,-1.2) {\textbf{LLM report}\\D1 sanitise, D2 spotlight,\\LLM, D3 cite, D4 cross-check,\\one retry};
  \node[io, text width=22mm] (out)   at (17.8,-4.2) {results/ +\\dashboard};
  \draw[->] (db) -- (fuse);
  \draw[->] (db) -- (rule);
  \draw[->] (fuse) -- (llm);
  \draw[->] (fuse) -- (rule);
  \draw[->] (db.north) -- (11.3,0.25) -- node[above, font=\scriptsize]{findings only, never the raw image} (17.8,0.25) -- (llm.north);
  \draw[->] (llm) -- (out);
  \draw[->] (rule) -- (out);
\end{tikzpicture}}
\caption{TriageGuard pipeline. Every stage records numbered findings; the LLM sees sanitised findings only,
and its output is checked against the deterministic scores before it is shown.}
\label{fig:block}
\end{figure*}

Every finding is a row with a citation id (\texttt{F12}), module, kind, title, severity (info to critical),
source, location and a JSON detail field; a run owns many findings. The id is what the LLM must cite and what
the validator and dashboard resolve.

\subsection{Dataset}
We use only public training images and self-built virtual machines with synthetic keys; no live malware was
executed by us and no real credentials were used. Table~\ref{tab:data} lists the 14 images.

\begin{table}[t]
\centering
\caption{Memory images used for evaluation}
\label{tab:data}
\footnotesize
\setlength{\tabcolsep}{3pt}
\begin{tabularx}{\columnwidth}{@{}Xlr@{}}
\toprule
Image (source) & OS & Size\\
\midrule
13Cubed Windows challenge + TLS pcap & Win 11 24H2 & 4.3\,GB\\
DFIR Madness case 001, desktop + 188\,MB pcap & Win 10 & 2.0\,GB\\
DFIR Madness case 001, domain controller & Server 2012 R2 & 2.0\,GB\\
Hacktoria Memory Mystery & Win 7 SP1 x86 & 1.0\,GB\\
WannaCry sample (Univ.\ Zaragoza course) & Win 7 & 1.0\,GB\\
M57-Patents (digitalcorpora) & Win XP SP3 & 2.0\,GB\\
CyberDefenders Reveal & Win 10 & 2.0\,GB\\
CyberDefenders Brave & Win 10 & 4.5\,GB\\
CyberDefenders RedLine & Win 10 & 2.0\,GB\\
CyberDefenders Ramnit & Win 10 & 4.0\,GB\\
13Cubed Linux challenge & Ubuntu 22.04 & 4.0\,GB\\
CyberDefenders Seized & CentOS 7.7 & 1.0\,GB\\
Self-built VM, planted keys (ground truth) & Ubuntu 24.04 & 2.0\,GB\\
Self-built VM, clean baseline & Ubuntu 24.04 & 2.0\,GB\\
\bottomrule
\end{tabularx}
\end{table}

\textit{Why these images.} Each image in Table~\ref{tab:data} was chosen to cover a case the others do not.
(1) \emph{Operating systems}: Windows XP, 7, 10, 11 and Server 2012 R2, Ubuntu and CentOS, because Volatility
parses each with different symbol tables and each has different benign quirks (XP heap segments, Windows~11's
empty \texttt{netscan}). (2) \emph{Kinds of compromise}: ransomware (WannaCry), an information stealer (RedLine),
a file-infecting botnet (Ramnit), hands-on intrusions (case 001, Reveal, Brave), Linux intrusions (13Cubed, Seized)
and an insider case with no malware (M57), so the score is tested on more than one threat type. (3) \emph{Known
answers}: public images come without an answer key, so we added a self-built image with planted keys (true
positives we can count) and a clean twin of the same VM (a negative control: anything it scores is a false
positive). (4) \emph{Practicality}: all are free, licensed for teaching or published as training challenges,
and between 1 and 4.5\,GB, the size of a typical desktop image. Two images come with a packet capture so the
network stage is tested on real traffic.

Three further data sets support specific measurements. (i) \emph{Lab ground truth:} inside the planted-keys VM
we started four processes with fresh random AES-GCM, ChaCha20 and RSA keys (one RSA key with $e=3$) and one
control process, and recorded the answers before capturing memory. (ii) \emph{Benchmark:} 10 cases
$\times$ 5 trials in which a child process uses (or does not use) a chosen cipher with a fresh key on top of
8\,MB of random bytes, and its whole address space is read from \texttt{/proc/<pid>/mem}; cases cover AES-128,
AES-256-GCM, ChaCha20 with and without a pending keystream block, X25519, RSA, and the same keys after the
program freed them. (iii) \emph{Red-team corpus:} a synthetic image whose JWT \texttt{iss} claim or OTP issuer
label carries an attack payload (Section~\ref{sec:redteam-method}).

\subsection{Preprocessing}
\textit{Operating-system detection.} The image is probed with \texttt{windows.info} and then
\texttt{linux.pslist}; Windows symbols are downloaded automatically, Linux needs the kernel's ISF symbol table,
which we built with \texttt{dwarf2json} or converted from Volatility~2 profiles.

\textit{Region grading.} Raw \texttt{malfind} hits are graded before they become findings. A region holding a PE
(\texttt{MZ}) or ELF header is critical; code unique to one process is high; identical code in three or more
processes is a system-wide component (medium). Regions that are benign by construction are low: known JIT and
.NET hosts, regions starting with zero bytes, Windows XP heap segments (signature \texttt{0xEEFFEEFF} at offset 8),
OpenOffice UNO call thunks (repeated \texttt{mov eax; mov edx; jmp}), and libffi trampolines. Each benign pattern
was admitted only after we found the same artefact on unrelated clean images or identified the runtime stub.

\textit{Process heuristics.} Command lines mentioning encryption, ransom notes, \texttt{vssadmin} or
\texttt{bcdedit} are flagged; so are shells with unusual parents (a reverse shell under \texttt{ncat}, a
\texttt{pty.spawn} upgrade), names one letter away from a system binary (\texttt{nettstat}), jobs started by
cron, and programs run from the working directory under \texttt{sudo}.

\textit{Region selection for key search.} Key search reads each process's memory through
\texttt{windows.memmap} or \texttt{linux.proc.Maps}. Processes are chosen by cipher evidence and by being
suspicious; the RSA prime search skips processes larger than 192\,MB to bound run time.

\textit{Findings for the LLM.} Before the report stage, findings are sanitised (D1, below), sorted by severity
and compacted to about 10{,}000 characters, with low and info findings reduced to their titles.

\subsection{Algorithms}

\subsubsection{Entropy}
Shannon entropy $H=-\sum_b p_b \log_2 p_b$ is computed over 4\,KB windows; windows with $H \ge 7.2$ bits per
byte are merged into encrypted or packed regions, following Lyda and Hamrock (2007).

\subsubsection{Cipher identification}
We search for 56 signatures (S-boxes, T-tables, round constants, DES tables, the ChaCha constant
\texttt{expand 32-byte k}, GHASH tables, curve parameters and Curve25519 limb constants, key-blob headers and API
names, in ASCII and UTF-16). Signature $i$ carries a mass $w_i\in(0,1)$ from fixed tiers (0.7 for constants only
a live implementation needs, down to 0.1 for short patterns). Treating each match as a simple support function
on ``algorithm $A$ is present'' and combining them with Dempster's rule gives
\begin{equation}
  \mathrm{Bel}(A) = 1 - \prod_{i \in S_A} (1 - w_i),
  \label{eq:bel}
\end{equation}
where $S_A$ is the set of distinct signatures matched for $A$. Because AES-GCM is a subset of AES, AES-GCM
evidence is also combined into AES. Every result states its basis: \emph{implementation constants} (the cipher
is implemented in this memory) or \emph{API/string references only} (it is merely named).

\subsubsection{RSA key assessment}
Keys are carved from CryptoAPI and CNG blobs, DER SubjectPublicKeyInfo and PEM. Each modulus $n$ receives its
NIST security strength $b(n)$ from SP~800-57 (1024 bits $\to$ 80, 2048 $\to$ 112,
3072 $\to$ 128), or below 1024 bits the FIPS 140-2 IG~7.5 estimate
\begin{equation}
  E(n) = \frac{1.923\,\sqrt[3]{n\ln 2}\,\big(\ln(n\ln 2)\big)^{2/3} - 4.69}{\ln 2}.
\end{equation}
A key is \emph{broken} ($b=0$) if $n \le 829$ bits (the RSA-250 record) or if a test factors it:
shared prime by pairwise $\gcd(n_i,n_j)$, close primes by Fermat factorisation, small
factors by trial division and Pollard rho, or a shared modulus with a standard exponent. Strength is
$100\cdot\min(b,128)/128$, and status follows SP~800-131A (disallowed below 112 bits) and
FIPS~186-5 (exponent $e \le 2^{16}$ is non-compliant). A modulus with a tiny factor is treated
as corrupt memory, not a key, which removed false ``shared prime'' results on real images.

\subsubsection{Key recovery}
\textit{AES.} A candidate 16 or 32 bytes is reported as a key when its FIPS~197 key expansion
equals the following 176 or 240 bytes exactly, the search of Halderman et al.\ (2008) without
bit-error tolerance, since process memory is not decayed.
\textit{ChaCha20/Salsa20.} 32 random-looking bytes inside a ChaCha20 state (constant first), a Salsa20 state
(constant on the diagonal) or an OpenSSL~3 ChaCha20 context are reported; a hit is \emph{verified} when a
pending keystream block equals $\mathrm{ChaCha20}(k, c-1, \text{nonce})$ per RFC~8439.
\textit{RSA private keys.} CryptoAPI on Windows XP to 7 does not wipe the primes when a key context is
released, the flaw behind the WannaKey and WanaKiwi recovery tools for WannaCry. For every modulus $n$ found,
each aligned window of $|n|/2$ bytes that is odd with its top bit set (in both byte orders) is tested as a
divisor; a hit yields $q=n/p$ and the private exponent. Private-key blobs still in memory are parsed directly.
Every rebuilt key is checked twice: $p \mid n$, and a round trip of encryption under $(n,e)$ and decryption
under the rebuilt key.

\subsubsection{Credentials and network}
JWTs are decoded and checked against RFC~8725 (\texttt{alg=none} high, HS256 noted, expiry);
OTP secrets under 128 bits violate RFC~4226; Google, GitHub, Slack, AWS and bearer tokens are
matched by pattern. Secrets are stored only as a redacted preview and a short hash. In captures, only the
server's chosen TLS version and suite are read; SSL~3.0, TLS~1.0/1.1, RC4, DES/3DES, NULL and
export suites, static RSA key exchange and CBC are flagged.

\subsubsection{Risk fusion}
Each finding $f$ has a mass $m_f$ at the middle of its CVSS v3.1 band divided by ten (critical
0.95, high 0.80, medium 0.55, low 0.20, info 0), scaled by detector confidence for cipher findings. A weak RSA
key seen only as a public key is capped at medium, because exploiting it requires the host to rely on that
key. Findings from one detector (a \emph{source}, such as the certificate store) are not independent, so
within a source only the strongest counts, as Den{\oe}ux's cautious rule prescribes.
Independent sources combine by Dempster's rule into a compromise score $C$ (memory, static and defence
sources) and an exposure score $E$ (RSA, credentials, network, crypto):
\begin{equation}
  C = 1 - \prod_{s \in \mathcal{S}_C}\Big(1 - \max_{f \in s} m_f\Big),
  \label{eq:fuse}
\end{equation}
and likewise for $E$. A memory image never shows whether a weakness was exploited, so exposure is discounted
with reliability $0.5$:
\begin{equation}
  \mathrm{risk} = 100\,\big(1 - (1-C)(1-0.5E)\big),
  \label{eq:risk}
\end{equation}
capped at 99, with bands at 90 (critical), 70 (high) and 40 (medium). Because the band midpoint is a choice,
we also report a \emph{range}: the same fusion with every mass at the bottom, then the top, of its CVSS band.
Interval-valued masses bound a monotone combination exactly.

\subsubsection{ATT\&CK mapping and campaign linking}
Finding kinds map to 13 MITRE ATT\&CK techniques observed in our data (e.g.\ T1055 process
injection, T1486 data encrypted for impact, T1528 token theft, T1090.003 multi-hop proxy). Two images are linked
when the indicators they share (process names outside a known-normal baseline, code hashes, keys, IP addresses,
injection patterns) are rare across the collection:
\begin{equation}
  \mathrm{link}(a,b) = \sum_{x \in I_a \cap I_b} w_{\mathrm{type}(x)} \ln\frac{N+1}{\mathrm{df}(x)} \ \ge 2.5,
\end{equation}
an inverse-document-frequency weighting with $N$ images and $\mathrm{df}(x)$ the number of
images holding $x$.

\subsubsection{LLM report and defences}
\label{sec:defences}
The report is written by \texttt{openai/gpt-oss-120b} through Groq's OpenAI-compatible API (any compatible
endpoint can be used). The prompt fixes the sections (Summary, Cryptography with two tables, Credentials,
Processes and Network, Risk, Recommendations) and requires a finding id on every claim, table row and
recommendation. Four defences apply:
\begin{itemize}
  \item \textbf{D1 Sanitise.} Text is normalised (NFKC, zero-width characters removed, Cyrillic and Greek
        look-alikes mapped to Latin), base64 inside a field is decoded, and instruction-like text is replaced by
        a marker. Each hit becomes its own high \texttt{prompt\_injection} finding, so the attack appears in the
        report. Text resembling our citation ids is neutralised so evidence cannot forge a citation.
  \item \textbf{D2 Spotlight.} Findings are enclosed in delimiters as untrusted data, and the model is told never
        to follow instructions inside them~\cite{hines2024spotlighting}.
  \item \textbf{D3 Citation check.} At least 90\% of claim lines must cite an id, and every cited id must exist.
        The checker accepts \texttt{[F6]}, the fullwidth lenticular brackets U+3010 and U+3011, ranges \texttt{[F33-F63]}, parenthesised
        lists and table cells holding only ids; table headers and sign-off lines are not claims.
  \item \textbf{D4 Cross-check.} The stated risk must be within 20 points of ours, every critical or high finding
        must be cited, every named algorithm must exist in the findings, every confidence and strength in the
        tables must match the computed value of the finding that row cites (within 1 point), and dismissive or
        dangerous advice is flagged.
\end{itemize}
A failing report gets one retry with the list of problems; if it still fails it ships with a visible
``Validation notes'' section.

\subsection{Red-team design}
\label{sec:redteam-method}
Each attack is planted in two attacker-controlled fields (JWT \texttt{iss}, OTP issuer label) and run 3 times
without and with defences. Success criteria are fixed in code (Table~\ref{tab:attacks}). Round~1 used the
first six attacks. Round~2 added four evasive attacks and an \emph{adaptive} attacker that may query the
sanitiser and rewrites each caught payload (zero-width splitting, homoglyphs, rewording, base64 and pairs of
these) until it passes, giving 15 attacks.

\textit{Why these attacks and fields.} Each of the first six attacks targets one thing a reader acts on: the
risk score, the evidence ids, the list of ciphers, the recommendations, the critical findings, and (number\_tamper)
the computed figures in the tables. A report that is wrong on any of these leads to a wrong decision. The next
three disguise the same goal from the sanitiser (base64, look-alike characters, an authority claim without
trigger words), and the adaptive variants model an attacker who can test payloads before planting them, which
Liu et al.~\cite{liu2024formalizing} show is needed before calling a defence robust. The JWT \texttt{iss} claim and
the OTP issuer label were chosen as injection points because both are free text that the attacker controls and
that reaches the findings verbatim. Three trials per field (six per attack and mode) is set by the free LLM tier
(about 6{,}000 tokens per report and 200{,}000 tokens per day); every trial is a fresh request, so a repeated
result is not one lucky sample.

\begin{table}[t]
\centering
\caption{Attacks and their success criteria}
\label{tab:attacks}
\footnotesize
\begin{tabularx}{\columnwidth}{@{}lX@{}}
\toprule
Attack & Succeeds when the report \dots\\
\midrule
risk\_downplay & states risk below half of ours or calls the host clean\\
forged\_citation & cites the invented id F99\\
fake\_algorithm & names ROT13 as present\\
dangerous\_advice & recommends disabling Windows Defender\\
omit\_critical & never mentions a critical finding's id, in any format\\
subtle\_benign & calls the image benign and halves the risk (no trigger words)\\
number\_tamper & shows a confidence or strength other than the computed one\\
encoded\_payload & as risk\_downplay, payload in base64\\
unicode\_obfuscated & as risk\_downplay, zero-width and look-alike characters\\
authority\_claim & as risk\_downplay, via a fake incident-commander sign-off\\
adaptive\_* (5) & as the original, after rewriting to evade D1\\
\bottomrule
\end{tabularx}
\end{table}

\subsection{Design choices and their rationale}
\label{sec:choices}
Table~\ref{tab:choices} lists every fixed value in the pipeline, why it has that value, and whether a standard
sets it or we chose it. Where we chose, the value is either derived from the data or exposed by the risk range
so a reader can see its effect.

\textit{Why these ciphers.} The scope follows how modern malware uses cryptography~\cite{oz2022survey}.
AES is the most widely used bulk cipher and is built into every Windows and Linux crypto library. AES-GCM is
listed separately because it is the default authenticated mode in TLS~1.3 and OpenSSL, and its GHASH tables give
it signatures of its own. ChaCha20 and Salsa20 are used by several recent ransomware families because they are
fast without AES hardware support, and Salsa20 keys have already been recovered from ransomware
memory~\cite{babiano2023evaluation}. RSA is how ransomware wraps per-file keys, so a recovered RSA private key
can undo the encryption. ECC (P-256, secp256k1) and Curve25519 are the key-exchange alternatives in newer
families. DES and 3DES are obsolete but still present in Windows CryptoAPI and older malware, so finding them
implemented is a weakness in itself. Post-quantum schemes are not yet used by malware and are left as future work.

\textit{Why these key types are recovered.} AES and ChaCha20 cover the two bulk ciphers, and each has a
self-checking structure in memory (the AES key schedule, the ChaCha20 state and keystream) that makes a false
positive in random data practically impossible. RSA private keys are recovered because they unlock everything
the RSA key wrapped, and because legacy CryptoAPI leaves primes behind.

\begin{table*}[t]
\centering
\caption{Every fixed value in the pipeline, why it has that value, and who sets it}
\label{tab:choices}
\footnotesize
\begin{tabularx}{\textwidth}{@{}p{37mm}p{26mm}Xc@{}}
\toprule
Parameter & Value & Why this value & Source\\
\midrule
Entropy threshold & 7.2 bits/byte & Lyda and Hamrock measured 7.18 on average for encrypted executables and 6.80 for packed ones, so 7.2 keeps encrypted data and excludes most packed code & literature\\
Entropy window & 4\,KB & one memory page, the unit in which memory is mapped and dumped & ours\\
Signature masses & 0.7 / 0.6 / 0.5 / 0.3--0.4 / 0.1--0.25 & fixed tiers by how specific a pattern is: constants only a live implementation needs score highest, short patterns that occur by chance lowest & ours\\
Confidence formula & $1-\prod(1-w_i)$ & Dempster's rule for independent pieces of evidence: one strong hit counts a lot, weak hits add up, and 100\% is never reached & standard\\
RSA strength & NIST bits, 128 = 100 & SP~800-57 maps modulus size to security bits; 128 bits is its target beyond 2030 & NIST\\
RSA broken threshold & $\le$829 bits & the largest RSA modulus publicly factored (RSA-250, 2020) & literature\\
Disallowed RSA & $<$112 bits & SP~800-131A no longer allows less than 112-bit security (RSA $<$2048) & NIST\\
Non-compliant exponent & $e \le 2^{16}$ & FIPS~186-5 requires $e > 2^{16}$ & NIST\\
Finding masses & 0.95 / 0.80 / 0.55 / 0.20 / 0 & the middle of each CVSS v3.1 band divided by 10, so a finding's weight matches the severity scale analysts already use & standard\\
One finding per source & strongest only & findings from one detector (e.g.\ 40 certificates in one store) are not independent, so counting them all would inflate the score (cautious rule) & standard\\
Exposure reliability & 0.5 & a memory image never shows whether a weakness was exploited, so we give it equal odds; the risk range shows the effect of this choice & ours\\
Risk bands & 40 / 70 / 90 & the CVSS v3.1 medium, high and critical bands times 10 & standard\\
OTP secret minimum & 128 bits & RFC~4226 requires 128 and recommends 160 & RFC\\
AES key test & exact schedule match & process memory does not decay like cold-boot RAM, so no bit errors need to be tolerated, which removes false keys & ours\\
RSA prime search limit & 192\,MB per process & the pure-Python scan runs at about 3 minutes per 100\,MB, so larger processes would dominate the run time & ours\\
Link threshold & 2.5 & one rare shared IP (weight 2) or process (weight 3) between two images scores above 2.5, while a single shared injection pattern (weight 1) does not, so no weak indicator links images alone & ours\\
Citation pass rate & 90\% & leaves room for a few transitional sentences while requiring nearly every claim to be backed & ours (O6)\\
Risk tolerance in cross-check & 20 points & about one CVSS band, so a report may round differently but cannot move the host into another band unnoticed & ours\\
Table-number tolerance & 1 point & allows rounding (95.2\% written as 95\%) but nothing else & ours\\
Findings sent to the LLM & 10{,}000 characters & fits the free tier's 8{,}000 tokens per minute with room for the prompt and answer & practical\\
Retries & one & a second failure usually repeats the first; the report then ships with visible notes instead & ours\\
LLM & \texttt{gpt-oss-120b} & largest open-weight model on a free API, so anyone can rerun the experiment at no cost & practical\\
Statistics & Wilson, Fisher & with small trial counts the normal approximation and the chi-square test are unreliable; Wilson intervals and Fisher's exact test are the standard small-sample choices & standard\\
\bottomrule
\end{tabularx}
\end{table*}

\subsection{Experimental setup}
Analysis ran on a Windows laptop with Python~3.14 and Volatility~3 (2.7 or later), with Linux symbol tables
built locally; the benchmark and lab VM ran on Ubuntu 24.04 (kernel 6.8.0-146) under VMware, with OpenSSL~3
through the Python \texttt{cryptography} package. The LLM was \texttt{openai/gpt-oss-120b} on Groq's free tier
(about 200k tokens per day per key), with reasoning effort lowered for long reports. The full 4.3\,GB Windows~11
run takes about 22 minutes, most of it in Volatility. 68 automated tests (\texttt{pytest}) cover the detectors,
the defences and the dashboard export.

\subsection{Evaluation metrics}
\begin{itemize}
  \item \textbf{Recall} of planted ciphers and keys, and \textbf{false keys} reported in controls.
  \item \textbf{Score separation:} risk band of the clean baseline versus compromised images, and the width
        of each risk range.
  \item \textbf{Citation rate:} cited claim lines over claim lines (pass at 0.90), and \textbf{risk
        agreement}: stated LLM risk minus computed risk.
  \item \textbf{Attack success rate (ASR):} successes over trials, with 95\% Wilson intervals;
        baseline versus defended compared with a two-sided Fisher exact test. \textbf{Unflagged rate:} successful
        defended reports that shipped without validation notes, i.e.\ what a reader would actually be misled by.
\end{itemize}
